\documentclass[twocolumn,10pt,a4paper]{article}
\usepackage[utf8]{inputenc}
\usepackage{amsmath, amssymb}
\usepackage{graphicx}
\usepackage{geometry}
\geometry{a4paper, margin=0.75in, columnsep=0.25in}
\usepackage{hyperref}
\usepackage{mathpazo} % Palatino font for a slightly more polished look
\usepackage{microtype}
\usepackage{float}
\usepackage{caption}
\usepackage{subcaption}

% Adjust title format to look like Arxiv
\makeatletter
\renewcommand{\maketitle}{\bgroup\setlength{\parindent}{0pt}
\begin{flushleft}
  {\Large\bfseries \@title}
  \vspace{1.5ex}
  
  {\normalsize \@author}
  \vspace{1.5ex}
  
  {\small \@date}
\end{flushleft}\egroup
}
\makeatother

\title{Notes on Femtosecond Pulse Generation:\\ Measurements and Optical Characterization}
\author{Gregorio Jaca and Peter Tallosy}
\date{\today}

\begin{document}

\twocolumn[
  \begin{@twocolumnfalse}
    \maketitle
    \vspace{-1em}
    \begin{abstract}
      \noindent These notes summarize the experimental characterization of a femtosecond pulsed laser system, serving as a two-part laboratory series. We detail the fundamental parameters of the pulse generator, spectral bandwidth measurements, and time-domain pulse width estimations assuming Fourier-transform limits. Furthermore, we discuss the role of saturable absorbers (SESAM) in mode-locking, the effects of dispersion from glass optical components, and the subsequent compression of pulses using an SF10 prism sequence aligned at Brewster's angle alongside Tellurium dioxide ($\text{TeO}_2$) crystals. Finally, we demonstrate the high-precision measurement of angular dispersion using spectrally resolved interferometry with a Michelson interferometer.
    \end{abstract}
    \vspace{0.5cm}
  \end{@twocolumnfalse}
]

\section{System Overview and Cavity Dynamics}
Our experimental setup is driven by a Continuous Wave (CW) pump laser operating at 20\,W, which is subsequently frequency-doubled via Second Harmonic Generation (SHG) down to 5.2\,W. In the SHG process, each photon doubles in energy, halving the original wavelength. This output pumps the Tsunami pulse generator, producing an output power of $P = 0.340$\,W in the red-infrared spectrum.

Unlike CW lasers, pulse operation requires the phase synchronization of numerous longitudinal modes, leading to a much broader spectrum. The formation of the pulse strictly relies on stimulated emission. Spontaneous emission (fluorescence), which carries random phase and frequency, does not contribute to laser action and is selectively suppressed by the resonant cavity to ensure a synchronized phase and fixed frequency.

The time between pulses (the round-trip time of the light inside the cavity) was measured to be $\Delta t = 12.14$\,ns. The cavity length $L$ can be calculated directly from this round-trip time:
\begin{equation}
    L = \frac{c \cdot \Delta t}{2} \approx \frac{(3 \times 10^8\,\text{m/s})(12.14\,\text{ns})}{2} \approx 1.82\,\text{m}.
\end{equation}

In any measurement chain, the slowest device determines the overall error. If the risetime for a detector is relatively higher than the pulse duration, the resulting measurement essentially resembles a triangle wave. For instance, the width of the pulse measured directly on the oscilloscope is around $2$\,ns, specifically because the oscilloscope is fundamentally a 500\,MHz class device, unable to truly resolve femtosecond features.

\section{Spectral Measurements and Transform Limits}
To accurately characterize the pulse width independently of electronic risetime, we used a spectrometer to measure intensity as a function of frequency. The initially mode-locked spectral profile displays a distinct bandwidth of $\Delta \lambda = 14.5$\,nm, bounded by $\lambda_1 = 778$\,nm and $\lambda_2 = 792.5$\,nm. 
In frequency space, these correspond to:
\begin{equation}
    f_1 = \frac{c}{\lambda_1} \approx 385.34\,\text{THz}, \quad f_2 = \frac{c}{\lambda_2} \approx 378.29\,\text{THz},
\end{equation}
yielding a spectral bandwidth of $\Delta f \approx 7.05$\,THz.

According to the time-bandwidth product, the pulse duration $\Delta t$ and spectral bandwidth $\Delta f$ obey $\Delta f \cdot \Delta t \ge C_B$, where $C_B$ depends on the temporal shape of the pulse. Assuming a Gaussian pulse shape, the constant is $C_B \approx 0.441$. Thus, the Fourier-transform-limited pulse duration is:
\begin{equation}
    \tau_{\text{pulse}} \approx \frac{0.441}{\Delta f} \approx 62.55\,\text{fs}.
\end{equation}
Our measured initial output pulse demonstrated an autocorrelation full-width at half-maximum (FWHM) of $61\,\text{fs}$, excellently matching the transform limit. A full spectrum comparison of the pulse profile before and after massive optical dispersion modifications is shown in Figure~\ref{fig:spectra}.

\begin{figure}[H]
     \centering
     \begin{subfigure}[b]{0.48\linewidth}
         \centering
         \includegraphics[width=\linewidth]{./figures/all_spectra.pdf}
         \caption{Normalized spectral profiles.}
         \label{fig:spectra_all}
     \end{subfigure}
     \hfill
     \begin{subfigure}[b]{0.48\linewidth}
         \centering
         \includegraphics[width=\linewidth]{./figures/FTS_from_time.pdf}
         \caption{Derived FTS spectrum.}
         \label{fig:fts_spectrum}
     \end{subfigure}
     \caption{Spectral characterization. (a) Displays the measured bandwidth evolution through different optical stages. (b) Shows the theoretical optical spectrum derived strictly from the time-domain pulse envelope via Fourier Transform Spectroscopy (FTS), compared with the measured data.}
     \label{fig:spectral_group}
\end{figure}

\section{Fourier-Domain Validation of the Transform Limit}
According to the Wiener-Khinchin theorem, the absolute intensity profile of the power spectrum can be directly tied to the temporal field envelope via the Fourier transform. We rigorously demonstrated this relationship computationally by utilizing the explicitly extracted initial optical spectrum dataset.

By systematically converting the recorded wavelength trace cleanly into the linear frequency domain, shifting the fundamental envelope center carrier frequency safely to baseband, assuming a strictly uniform flat-phase representation perfectly across the active bandwidth (the rigorous formal definition of the Fourier-transform limit), and heavily zero-padding the subsequent field envelope block ensuring maximum continuous time-domain evaluation interpolation, an Inverse Fast Fourier Transform (IFFT) natively yields the predicted optimal shortest intensity profile curve boundary.

Plotting the strictly discrete intensity auto-correlation calculation tracking tracking this synthetic IFFT profile ($I_{\text{ac}}(\tau) = \int I(t) I(t-\tau) dt$) solidly directly alongside the purely unfiltered raw experimental hardware autocorrelation trace, as rendered comprehensively cleanly in Figure~\ref{fig:fourier_validation}, definitively actively proves that our continuously tracked $\sim 61\,\text{fs}$ initial stage pulse genuinely operates incredibly flawlessly right at the ideal theoretical physical transform barrier constraint bounds, firmly inherently establishing practically zero measurable internal residual pre-chirp.

\begin{figure}[H]
    \centering
    \includegraphics[width=\linewidth]{./figures/FT_validation.pdf}
    \caption{Continuous robust overlay plot strictly demonstrating the directly measured raw lab autocorrelation traces precisely stacked exclusively against the purely mathematically derived theoretical Transform-Limited (TL) bounds curve computed natively solely via discrete IFFT bulk array manipulation of the isolated optical frequency spectrum, permanently verifying continuous pulse quality limitations cleanly.}
    \label{fig:fourier_validation}
\end{figure}

\section{Mode-locking and Autocorrelation}
Our system uses a Semiconductor Saturable Absorber Mirror (SESAM) to passively mode-lock the laser and force pulsed operation. The SESAM effectively acts as an intensity-dependent filter; it lets high-intensity coherent pulses pass unharmed while strongly absorbing low-intensity continuous noise, continually polishing the cavity contents without introducing active spontaneous amplification.

We evaluated the pulse shape using an autocorrelator. Mathematically, the pulse envelope formation modeled with the saturable absorber leads to a differential equation whose complex solution intrinsically favors a hyperbolic secant squared ($\text{sech}^2$) temporal profile, though Gaussian profiles are also commonly approximated. We extracted both Gaussian and $\text{sech}^2$ fits from our experimental autocorrelator data continuously throughout to quantify our adjustments. As demonstrated in Figure~\ref{fig:initial_fits}, the initial pulse profile is better represented by the $\text{sech}^2$ model, which yielded a higher $R^2$ value and pulse duration more consistent with our spectral bandwidth.

\begin{figure}[H]
     \centering
     \begin{subfigure}[b]{0.48\linewidth}
         \centering
         \includegraphics[width=\linewidth]{./figures/TBP_autocorr1_61fs_leftpulse_sech2.pdf}
         \caption{$\text{sech}^2$ Fit.}
         \label{fig:initial_sech}
     \end{subfigure}
     \hfill
     \begin{subfigure}[b]{0.48\linewidth}
         \centering
         \includegraphics[width=\linewidth]{./figures/TBP_autocorr1_68fs_leftpulse_gauss.pdf}
         \caption{Gaussian Fit.}
         \label{fig:initial_gauss}
     \end{subfigure}
     \caption{Initial pulse autocorrelation fits for the oscillator output (Left pulse). The $\text{sech}^2$ profile demonstrates superior fit quality for the mode-locked output.}
     \label{fig:initial_fits}
\end{figure}

\section{Dispersion and Pulse Compression}

\subsection{Material Dispersion}
Optical amplifiers reduce the spectrum bandwidth, making the pulse longer in time. Furthermore, placing refractive glass in the transmissive beam path directly causes Group Delay Dispersion (GDD) or ``chirp", wherein the temporal width physically stretches because dispersive spectral components travel at minutely distinct group velocities.

We cleanly demonstrated this dispersion by reflecting the laser pulse using a 1\,cm thick plain glass reflector piece. We successfully isolated two reflections: one of them reflects rigidly from the front surface and does not go through the bulk glass (labeled ``Left''), while the other reflection traverses the internal glass width twice at an angle (labeled ``Right''). The ``Right'' reflection accrues measurable material dispersion, inherently becoming wider in time. As shown in Figure~\ref{fig:leftright}, the Right pulse effectively broadens to $\sim 119\,\text{fs}$ compared to the undispersed initial $\sim 61\,\text{fs}$ Left pulse, strictly validating the broadening effect.

\begin{figure}[H]
     \centering
     \begin{subfigure}[b]{0.48\linewidth}
         \centering
         \includegraphics[width=\linewidth]{./figures/TBP_autocorr1_61fs_leftpulse_sech2.pdf}
         \caption{Left Pulse (Undispersed).}
         \label{fig:left_pulse}
     \end{subfigure}
     \hfill
     \begin{subfigure}[b]{0.48\linewidth}
         \centering
         \includegraphics[width=\linewidth]{./figures/TBP_autocorr1_119fs_rightpulse_sech2.pdf}
         \caption{Right Pulse (Dispersed).}
         \label{fig:right_pulse}
     \end{subfigure}
     \caption{Direct comparison of the autocorrelation traces for the front surface reflection (Left) and the back surface reflection (Right). The Right pulse exhibits significant temporal broadening ($\sim 115\,\text{fs}$) compared to the initial pulse ($\sim 61\,\text{fs}$) due to travel through 1\,cm of extra glass material.}
     \label{fig:leftright_individual}
\end{figure}

\subsection{Pulse Compression Methodologies}
To precisely compress and regulate the pulse time, an SF10 high-dispersion glass prism cascade configuration is applied as a temporal compressor. Crucially, the constituent prisms are actively aligned carefully at the \textit{Brewster angle} to strictly minimize ordinary reflection losses for the incoming p-polarized laser beam profile. The prism arrangement intentionally generates a net negative GDD path that selectively delays longer wavelengths behind the shorter ones, perfectly counter-balancing and compensating for the unwanted positive dispersion of ordinary materials previously inserted.

For rigorous macroscopic testing, we also deliberately inserted Tellurium dioxide ($\text{TeO}_2$) crystals into the path. As shown in Figure~\ref{fig:teo2_comp}, the pulse was significantly broadened but partially compensated. The final optimized system performance, including all dispersive elements and the SF10 prism compressor, is detailed in Figure~\ref{fig:final_comp}, where the pulse width was restored to its near-transform-limited state.

\begin{figure}[H]
     \centering
     \begin{subfigure}[b]{0.48\linewidth}
         \centering
         \includegraphics[width=\linewidth]{./figures/TBP_autocorr_comp_TeO2_extended_65fs_sech2.pdf}
         \caption{$\text{sech}^2$ Fit.}
         \label{fig:teo2_sech}
     \end{subfigure}
     \hfill
     \begin{subfigure}[b]{0.48\linewidth}
         \centering
         \includegraphics[width=\linewidth]{./figures/TBP_autocorr_comp_TeO2_extended_72fs_gauss.pdf}
         \caption{Gaussian Fit.}
         \label{fig:teo2_gauss}
     \end{subfigure}
     \caption{Pulse characterization after 4\,cm of $\text{TeO}_2$ insertion. The $R^2$ values indicate that even under massive dispersion, the temporal profile remains consistent with our fitting models.}
     \label{fig:teo2_comp}
\end{figure}

\begin{figure}[H]
     \centering
     \begin{subfigure}[b]{0.48\linewidth}
         \centering
         \includegraphics[width=\linewidth]{./figures/TBP_autocorr_comp_shifted_prism_TeO2_glass_64fs_sech2.pdf}
         \caption{$\text{sech}^2$ Fit.}
         \label{fig:final_sech}
     \end{subfigure}
     \hfill
     \begin{subfigure}[b]{0.48\linewidth}
         \centering
         \includegraphics[width=\linewidth]{./figures/TBP_autocorr_comp_shifted_prism_TeO2_glass_72fs_gauss.pdf}
         \caption{Gaussian Fit.}
         \label{fig:final_gauss}
     \end{subfigure}
     \caption{Final optimized pulse compression results. The system successfully compensates for the cumulative dispersion of glass blocks and $\text{TeO}_2$ crystals, achieving a $\sim 64\,\text{fs}$ FWHM.}
     \label{fig:final_comp}
\end{figure}


\section{Interferometric Angular Dispersion Measurement}

\subsection{Spectrally Resolved Interferometry}
For the secondary investigation occasion, we pivoted to strictly evaluate the angular dispersion parameters generated inherently by wedge glass cuts or dispersive optical testing samples. An advanced experimental standard method for calculating and recovering this relies upon continuous spectrally resolved interferometry. Using a classically configured Michelson layout, the primary incoming test pulse is evenly divided via a generic beam-splitter block into two corresponding arms. One simple path acts strictly as a static baseline reference arm, while the other extended arm necessarily crosses through the specific dispersive optical element (in this immediate demonstration, a standard $60^\circ$ testing prism block).

Due directly to slightly mismatched baseline arm lengths and the introduced differential angular spectrum spread dispersed actively by the optical sample, the eventually recombined exiting beams intrinsically form a deeply distinct cross-plane spatial and spectral interference mesh array pattern (an observable stationary 2D field interferogram map). The visible relative macroscopic orientation shifts and linear tracking density layout geometry of these exact interference field banding fringes mathematically correlates with extremely precise linearity directly back to the physical angular separation spanning the spectrum.

\subsection{Data Extraction and Verification}
Figure~\ref{fig:interferogram} displays standard output arrays of the mapped 2D sensor interferogram frames precisely documented via static focal tracking optics at the output tail of the Michelson block layout. This recorded visualization neatly spans a continuous horizontal baseline mapped dynamically representing the input laser wavelength continuous sweep, combined securely with an active independent lateral vertical axis denoting basic independent beam spatial dimensions. The explicit planar geometric angular tilt shift explicitly shown strictly within the interference banding fields corresponds flawlessly to mapping proportional active angular dispersion limits directly transferred from the mounted sample.

\begin{figure}[H]
     \centering
     \begin{subfigure}[b]{0.48\linewidth}
         \centering
         \includegraphics[width=\linewidth]{./figures/michelson_interferogram.pdf}
         \caption{2D Sensor Interferogram.}
         \label{fig:interferogram_2d}
     \end{subfigure}
     \hfill
     \begin{subfigure}[b]{0.48\linewidth}
         \centering
         \includegraphics[width=\linewidth]{./figures/michelson_spectrum.pdf}
         \caption{Extracted 1D Spectral Profile.}
         \label{fig:michelson_spectrum}
     \end{subfigure}
     \caption{Direct sensor capture and mathematical extraction from the Michelson interferometer. (a) Displays the 2D interferogram field. (b) Shows the high-fidelity 1D spectral profile recovered by vertically integrating the interferogram matrix, which effectively averages out the interference fringes and noise to isolate the spectral envelope.}
     \label{fig:interferometry_group}
\end{figure}

Extracting calculations actively straight from observing and manually interrogating the structured interferogram imagery frame boundaries, our core angular dispersion property ($\Delta \gamma$) is explicitly verified effectively strictly analyzing internal fringe spatial coordinate scaling mappings. Operating fundamentally based on physical experimental cross-references, operating utilizing a total horizontal span wavelength capture envelope size spanning practically exactly $\Delta \lambda = 25$\,nm, locked onto an independently verified active vertical matrix digital array pixel calibration density scaling of $20\,\mu\text{m/pixel}$, combined simultaneously with a manually resolved and locked primary measured spatial coordinate consecutive lateral fringe peak-to-peak field distance $\Delta x = 240\,\mu\text{m}$, the physically derived raw spatial angular optical tilt spread mapped independently purely from inspecting consecutive fringes rigidly calculates practically effortlessly into resolving out to roughly an incredibly small scale spread factor approximating exactly out to $2.08 \times 10^{-3}$\,rad overall across the isolated test envelope. Actively mathematically standardizing this directly corresponds to resolving out a finalized clean absolute unit angular dispersion performance value equating to:
\begin{equation}
    \frac{d\gamma}{d\lambda} \approx \frac{2.08\,\text{mrad}}{25\,\text{nm}} \approx 83\,\mu\text{rad/nm}.
\end{equation}

Ultimately, operating completely uncoupled, this advanced experimental arrangement accurately underscores and actively highlights the uniquely astonishing precision threshold envelope sensitivities practically intrinsically inherent to utilizing true active continuous cross-plane spectrally resolved optical interferometry testing topologies over utilizing generic single-point position tracking diode sensors uniformly, continuously actively permitting and establishing clean optical characterizations mapping absolutely stable profiles natively extracting remarkably subtle physical properties from even barely mildly actively dispersive experimental material samples easily tracking down reliably efficiently cleanly handling extreme dimensional thresholds.

\end{document}
